\documentclass[11pt]{article}

\usepackage[margin=1in]{geometry}
\usepackage{amsmath,amssymb,mathtools}
\usepackage{bm}
\usepackage{hyperref}
\usepackage{xcolor}
\usepackage{listings}

\lstdefinelanguage{Julia}{
  morekeywords={function,end,for,in,if,else,elseif,return,using,const,Dict,Vector,Int,Float64,String,Any,true,false,do},
  sensitive=true,
  morecomment=[l]\#,
  morestring=[b]",
}

\lstset{
  language=Julia,
  basicstyle=\ttfamily\small,
  keywordstyle=\color{blue!70!black},
  commentstyle=\color{green!40!black},
  stringstyle=\color{red!50!black},
  showstringspaces=false,
  breaklines=true,
  frame=single,
  columns=fullflexible,
}

\newcommand{\ket}[1]{\left|#1\right\rangle}
\newcommand{\bra}[1]{\left\langle#1\right|}
\newcommand{\braket}[2]{\left\langle#1\,\middle|\,#2\right\rangle}
\newcommand{\V}{\mathcal V}
\newcommand{\FQ}{F_Q}
\newcommand{\FEP}{F_{\rm EP}}
\newcommand{\dd}{\delta}
\newcommand{\etaK}[1]{\eta^{\star,\mathrm{conn}}_{\le #1}}
\newcommand{\Tr}{\mathrm{Tr}}

\title{Operator-Resolved QFI Response and Matched-Filter Computation of \texorpdfstring{$\eta_{\le k}$}{eta<=k}}
\author{QFI chain notes}
\date{\today}

\begin{document}
\maketitle

\section{Physics Question}

The sharp question is not whether a single scalar ``detects
entanglement.''  The sharp question is:
\begin{quote}
  Given a state $\ket{\psi}$ and a generator $G$, through which
  observable algebra is the QFI response
  $\ket{\dd G}=(G-\langle G\rangle)\ket{\psi}$ visible?
\end{quote}
The total response strength is measured by
\[
  \FQ=4\braket{\dd G}{\dd G}.
\]
The matched-filter ratio
\[
  \eta^\star_{\V}
  =
  \frac{
    \|P_{\V\ket{\psi}}\ket{\dd G}\|^2
  }{
    \|\ket{\dd G}\|^2
  }
\]
then asks what fraction of that response can be read out by observables
inside a specified real Hermitian algebra $\V$.

Thus $\eta^\star_{\V}$ is best understood as an
\emph{operator-resolved visibility of the QFI response}.  It is not a
universal entanglement monotone and should not be interpreted as a
sure-fire diagnostic of long-range entanglement by itself.  The choice
of $\V$ is part of the physical question.

\paragraph{Example: Haldane-chain warning.}
In a spin-1 Haldane phase, local two-point correlations are short-ranged,
but the entanglement/correlation structure is naturally exposed by
nonlocal string operators.  If one restricts $\V$ to ordinary local
windows, $\eta^\star_{\V}$ may remain small until large windows are
included.  That does \emph{not} mean the Haldane phase is long-range
entangled in the GHZ sense.  It means the response to the chosen $G$ is
not efficiently visible in the local-window algebra.  If $\V$ is enlarged
to include the appropriate string algebra, the same short-range
entangled SPT state can have a large $\eta^\star_{\V}$.

This is the correct interpretation:
\[
  \boxed{
  \eta^\star_{\V}
  \text{ measures visibility through } \V,
  \text{ not entanglement itself.}
  }
\]
Different choices of $\V$ expose different channels:
\begin{itemize}
  \item local-window algebras diagnose local recoverability of the
  response;
  \item string algebras diagnose hidden or SPT response channels;
  \item scale-dependent families diagnose critical delocalization of the
  response.
\end{itemize}

\section{Numerical Purpose}

The important numerical observables for the corrected period-4
$J_1$--$J_2$--$J_3$ chain are
\begin{equation}
  \FQ[\psi;G]
  \qquad\text{and}\qquad
  \etaK{k}
  =
  \frac{\FEP[O^\star_{\le k}]}{\FQ},
\end{equation}
where $G=S^z_\pi$ is the staggered magnetization and
$O^\star_{\le k}$ is the optimal matched-filter observable inside the
real Hermitian subspace of connected Pauli windows of size at most $k$.
For this specific paper figure, $\V$ is the connected-window algebra,
so $\eta_{\le k}$ means ``local-window visibility of the staggered QFI
response.''
For the present path scan we need $k=2,4$ at
\[
  N=16,32,64
\]
along
\[
  D\to C:\quad (J_1,J_2,J_3)=(1,t,0),
\]
\[
  C\to U:\quad (J_1,J_2,J_3)=(1,1,t),
\]
\[
  U\to D:\quad (J_1,J_2,J_3)=(1,t,t).
\]
Here $J_1,J_2,J_3$ label the period-4 nearest-neighbor bond pattern
\[
  J_1,J_2,J_1,J_3,J_1,J_2,J_1,J_3,\ldots,
\]
not first-, second-, and third-neighbor couplings.

\section{Definitions}

Let $\ket{\psi}$ be a normalized DMRG ground state.  The generator is
\begin{equation}
  G
  =
  S^z_\pi
  =
  \sum_{j=1}^{N}(-1)^j S^z_j .
\end{equation}
The QFI response vector is
\begin{equation}
  \ket{\dd G}
  =
  \left(G-\langle G\rangle\right)\ket{\psi},
  \qquad
  \langle G\rangle=\bra{\psi}G\ket{\psi}.
\end{equation}
The pure-state QFI is
\begin{equation}
  \FQ
  =
  4\braket{\dd G}{\dd G}.
\end{equation}
In the ratio $\eta$, the overall normalization of $G$ cancels; the code
therefore only needs a consistent $G$ convention between numerator and
denominator.

\section{How \texorpdfstring{$F_Q$}{FQ} is Computed}

For a pure state,
\begin{equation}
  \FQ[\psi;G]
  =
  4\,\mathrm{Var}_{\psi}(G)
  =
  4
  \left(
    \langle G^2\rangle-\langle G\rangle^2
  \right).
  \label{eq:FQ_var}
\end{equation}
For the staggered generator
\[
  G=S^z_\pi=\frac{1}{\sqrt N}\sum_{j=1}^{N}(-1)^j S^z_j
\]
used in the production observable files, this becomes
\begin{align}
  \FQ
  &=
  4
  \left[
    \frac1N
    \sum_{j,k=1}^{N}
    (-1)^{j+k}
    \left(
      \langle S^z_j S^z_k\rangle
      -
      \langle S^z_j\rangle\langle S^z_k\rangle
    \right)
  \right].
  \label{eq:FQ_corr}
\end{align}
This is the formula used in the DMRG observable JSONs.
The factor $1/\sqrt N$ makes $G$ an intensive collective mode, so
$\FQ=O(1)$ for the dimer and cluster product states and grows only
logarithmically at the uniform Heisenberg point.

\subsection{Correlation-matrix implementation}

The code first computes the one-point magnetization
\begin{equation}
  m_j^z=\langle S^z_j\rangle
\end{equation}
and the two-point correlator
\begin{equation}
  C^{zz}_{jk}=\langle S^z_jS^z_k\rangle .
\end{equation}
For $j\ne k$, $C^{zz}_{jk}$ is obtained from ITensor's
\texttt{correlation\_matrix}; for $j=k$ the spin-$1/2$ identity gives
\begin{equation}
  \langle (S^z_j)^2\rangle=\frac14.
\end{equation}
The production routine in \texttt{observables.jl} implements
\eqref{eq:FQ_corr} as
\begin{lstlisting}
# F_Q (probe G = N^{-1/2} sum_j (-1)^j S^z_j; F_Q = 4 Var(G)).
FQ = 0.0
for j in 1:N, k in 1:N
    sgn = ((-1)^(j+k))
    czz = j < k ? real(two_pt[3, 3, j, k]) :
          j > k ? real(two_pt[3, 3, k, j]) :
          0.25
    FQ += 4.0/N * sgn * (czz - one_pt[3, j] * one_pt[3, k])
end
\end{lstlisting}
This is equivalent to building the MPO for $G$, applying it to the MPS,
and evaluating $4\braket{\dd G}{\dd G}$.

\subsection{MPO implementation}

The eta postprocessor uses the vector form directly.  It constructs
\[
  G=\sum_j(-1)^jS^z_j
\]
as an MPO, applies it to the MPS, and centers the result:
\begin{lstlisting}
Gpsi = apply(G, psi; cutoff=1e-14)
target = Gpsi - inner(psi, Gpsi) * psi
\end{lstlisting}
Then
\begin{equation}
  \braket{\mathrm{target}}{\mathrm{target}}
  =
  \braket{\dd G}{\dd G}
\end{equation}
is the denominator used in $\eta$.  The extra normalization by
$1/\sqrt N$ is unnecessary inside $\eta$, because both numerator and
denominator scale in the same way under $G\mapsto cG$:
\[
  \ket{\dd G}\mapsto c\ket{\dd G},
  \qquad
  \eta\mapsto \eta.
\]

\subsection{Why both \texorpdfstring{$F_Q$}{FQ} methods agree}

There are two equivalent computational routes:
\begin{enumerate}
  \item compute all $S^zS^z$ connected correlations and evaluate
  \eqref{eq:FQ_corr};
  \item build the MPS vector $\ket{\dd G}$ and evaluate
  $4\braket{\dd G}{\dd G}$.
\end{enumerate}
They are the same because
\begin{align}
  \braket{\dd G}{\dd G}
  &=
  \bra{\psi}
  \left(G-\langle G\rangle\right)^2
  \ket{\psi} \\
  &=
  \langle G^2\rangle-\langle G\rangle^2 \\
  &=
  \frac1N
  \sum_{j,k}
  (-1)^{j+k}
  \left(
    \langle S^z_jS^z_k\rangle
    -
    \langle S^z_j\rangle\langle S^z_k\rangle
  \right).
\end{align}
The JSON files store the first route because it is also useful as a
diagnostic of the raw spin correlations.  The eta postprocessor uses the
second route because it needs $\ket{\dd G}$ as an actual MPS target
vector for the matched-filter projection.

\section{Matched Filter on an Operator Subspace}

Let
\[
  \V_{\le k}^{\rm conn}
  =
  \mathrm{span}_{\mathbb R}\{O_a\}
\]
be the real vector space spanned by Hermitian Pauli strings supported on
connected windows of length at most $k$.
For a candidate observable
\[
  O(\bm w)=\sum_a w_a O_a,
\]
the error-propagation Fisher information is
\begin{equation}
  \FEP[O(\bm w)]
  =
  \frac{
    \left|-i\langle[O(\bm w),G]\rangle\right|^2
  }{
    \mathrm{Var}(O(\bm w))
  }.
\end{equation}
Define
\begin{equation}
  s_a
  =
  -i\langle[O_a,G]\rangle,
  \qquad
  \Sigma_{ab}
  =
  \frac12
  \left\langle
    \{\dd O_a,\dd O_b\}
  \right\rangle,
\end{equation}
where
\[
  \dd O_a=O_a-\langle O_a\rangle .
\]
Then
\begin{equation}
  \FEP[O(\bm w)]
  =
  \frac{(\bm w^T\bm s)^2}{\bm w^T\Sigma\bm w}.
\end{equation}

\subsection{Why the maximizer is the Wiener solution}

The optimization problem is therefore
\begin{equation}
  \max_{\bm w\ne 0}
  R(\bm w),
  \qquad
  R(\bm w)
  =
  \frac{(\bm w^T\bm s)^2}{\bm w^T\Sigma\bm w}.
  \label{eq:rayleigh}
\end{equation}
This is a generalized Rayleigh quotient.  The numerator is the squared
linear signal extracted by the observable $O(\bm w)$, while the
denominator is its variance/noise.  Thus the problem is exactly the
same algebraic problem as a matched filter or Wiener filter: choose
weights $\bm w$ that maximize signal-to-noise.

First assume that $\Sigma$ is strictly positive definite.  The quotient
is invariant under rescaling $\bm w\mapsto c\bm w$, so we may impose the
normalization constraint
\[
  \bm w^T\Sigma\bm w=1
\]
and maximize
\[
  (\bm w^T\bm s)^2.
\]
Equivalently, maximize $\bm w^T\bm s$ with the same constraint; the
overall sign of $\bm w$ is irrelevant because $\FEP$ squares the signal.
Introduce a Lagrange multiplier:
\begin{equation}
  \mathcal L(\bm w,\lambda)
  =
  \bm w^T\bm s
  -
  \frac{\lambda}{2}
  \left(
    \bm w^T\Sigma\bm w-1
  \right).
\end{equation}
Stationarity gives
\begin{equation}
  \frac{\partial\mathcal L}{\partial\bm w}
  =
  \bm s-\lambda\Sigma\bm w
  =
  0,
\end{equation}
so
\begin{equation}
  \Sigma\bm w
  =
  \lambda^{-1}\bm s.
\end{equation}
The overall scalar is immaterial, since $R(\bm w)$ is scale invariant.
Therefore the optimizing direction is
\begin{equation}
  \bm w^\star \propto \Sigma^{-1}\bm s.
\end{equation}
Choosing the natural normalization gives
\begin{equation}
  \bm w^\star=\Sigma^{-1}\bm s.
\end{equation}
Substituting back into \eqref{eq:rayleigh},
\begin{align}
  R(\bm w^\star)
  &=
  \frac{
    \left[
      (\Sigma^{-1}\bm s)^T\bm s
    \right]^2
  }{
    (\Sigma^{-1}\bm s)^T\Sigma(\Sigma^{-1}\bm s)
  } \\
  &=
  \frac{
    \left[
      \bm s^T\Sigma^{-1}\bm s
    \right]^2
  }{
    \bm s^T\Sigma^{-1}\bm s
  } \\
  &=
  \bm s^T\Sigma^{-1}\bm s.
\end{align}

Equivalently, the same result follows immediately from
Cauchy--Schwarz.  Write
\[
  \bm u=\Sigma^{1/2}\bm w,
  \qquad
  \bm v=\Sigma^{-1/2}\bm s.
\]
Then
\[
  \bm w^T\bm s
  =
  \bm u^T\bm v,
  \qquad
  \bm w^T\Sigma\bm w
  =
  \bm u^T\bm u.
\]
Therefore
\begin{equation}
  R(\bm w)
  =
  \frac{(\bm u^T\bm v)^2}{\bm u^T\bm u}
  \le
  \bm v^T\bm v
  =
  \bm s^T\Sigma^{-1}\bm s,
\end{equation}
with equality when
\[
  \bm u\propto \bm v,
  \qquad\text{i.e.}\qquad
  \Sigma^{1/2}\bm w
  \propto
  \Sigma^{-1/2}\bm s,
\]
again giving
\[
  \bm w^\star\propto\Sigma^{-1}\bm s.
\]

\subsection{What \texorpdfstring{$\Sigma^+$}{Sigma+} means}

In practice, $\Sigma$ need not be strictly invertible.  This is not a
numerical accident: the operator list often contains redundant
directions.  For example, two different Pauli strings may create the
same vector after acting on $\ket{\psi}$, or a string may annihilate the
state entirely.  Such directions have zero variance and therefore lie
in the null space of $\Sigma$.

The symbol
\[
  \Sigma^+
\]
means the Moore--Penrose pseudoinverse.  If
\begin{equation}
  \Sigma
  =
  V\Lambda V^T
  =
  \sum_{\alpha}
  \lambda_\alpha
  \bm v_\alpha\bm v_\alpha^T
\end{equation}
with $\lambda_\alpha\ge0$, then
\begin{equation}
  \boxed{
  \Sigma^+
  =
  \sum_{\alpha:\lambda_\alpha>0}
  \lambda_\alpha^{-1}
  \bm v_\alpha\bm v_\alpha^T.
  }
  \label{eq:pseudoinverse}
\end{equation}
Zero-variance directions are discarded rather than inverted.

With a singular covariance matrix, decompose the operator coefficient
vector into range and null components,
\[
  \bm w=\bm w_{\rm ran}+\bm w_{\rm null},
  \qquad
  \Sigma \bm w_{\rm null}=0.
\]
The variance depends only on $\bm w_{\rm ran}$.  A physically sensible
observable cannot obtain finite signal from a strictly zero-variance
direction; equivalently, the signal vector $\bm s$ lies in the range of
$\Sigma$ for exact arithmetic.  Numerically, we project onto the
positive-eigenvalue subspace.  The optimizing minimum-norm coefficient
vector is then
\begin{equation}
  \bm w^\star
  =
  \Sigma^+\bm s,
\end{equation}
and the optimal extracted Fisher information is
\begin{equation}
  \boxed{
  \eta^\star_{\V}
  =
  \frac{\bm s^T\Sigma^+\bm s}{\FQ}.
  }
\end{equation}

This is what the code implements.  In the notation of the MPS Gram
matrix below, $\Sigma$ is represented by the real symmetric Gram matrix
$K$, and $K^+$ is applied by diagonalizing $K$, discarding eigenvalues
below a tolerance, and dividing only by the retained eigenvalues.

\section{Equivalent Projection Formula}

For implementation, it is more stable to use the Hilbert-space
projection form.  Define
\begin{equation}
  \ket{\phi_a}
  =
  \dd O_a\ket{\psi}
  =
  \left(O_a-\langle O_a\rangle\right)\ket{\psi}.
\end{equation}
Then
\begin{equation}
  \eta^\star_{\V}
  =
  \frac{
    \left\|
      P_{\mathrm{span}_{\mathbb R}\{\ket{\phi_a}\}}
      \ket{\dd G}
    \right\|^2
  }{
    \braket{\dd G}{\dd G}
  }.
\end{equation}

For a real Hermitian operator subspace, the ED code uses the real-form
embedding
\begin{equation}
  \Phi_{\rm real}
  =
  \begin{bmatrix}
    \mathrm{Re}\,\phi_1 & \cdots & \mathrm{Re}\,\phi_M\\
    \mathrm{Im}\,\phi_1 & \cdots & \mathrm{Im}\,\phi_M
  \end{bmatrix},
  \qquad
  t
  =
  \begin{bmatrix}
    \mathrm{Im}\,\dd G\\
    -\mathrm{Re}\,\dd G
  \end{bmatrix}.
\end{equation}
Then
\begin{equation}
  \eta^\star_{\V}
  =
  \frac{\|P_{\mathrm{col}(\Phi_{\rm real})}t\|^2}{\|t\|^2}.
\end{equation}
The MPS code computes this same projection without expanding the full
$2^N$ vector.

\section{Symmetry Reduction: Odd-\texorpdfstring{$Y$}{Y} Sector}

The DMRG states in this campaign are real and live in the
$S^z_{\rm tot}=0$ sector.  The generator $G=S^z_\pi$ is real and
diagonal, so
\[
  \ket{\dd G}
\]
is a real vector.  In the real-form target
\[
  t
  =
  \begin{bmatrix}
    \mathrm{Im}\,\dd G\\
    -\mathrm{Re}\,\dd G
  \end{bmatrix},
\]
the first block vanishes and the second block is nonzero.

A Hermitian Pauli string with an even number of $Y$ operators maps a real
state to a real vector.  Its real-form column has support in the first
block:
\[
  \begin{bmatrix}
    \mathrm{Re}\,\phi_a\\
    0
  \end{bmatrix}.
\]
It therefore has zero direct overlap with $t$.

A Hermitian Pauli string with an odd number of $Y$ operators maps a real
state to an imaginary vector.  Its real-form column has support in the
second block:
\[
  \begin{bmatrix}
    0\\
    \mathrm{Im}\,\phi_a
  \end{bmatrix},
\]
and can carry matched-filter signal.

The second exact filter comes from the unbroken $U(1)$ spin symmetry.
The DMRG state lies in the $S^z_{\rm tot}=0$ sector, and
$G=S^z_\pi$ commutes with $S^z_{\rm tot}$.  Therefore
\[
  \ket{\dd G}
\]
also lies in the $S^z_{\rm tot}=0$ sector.  Any operator string whose
ladder-operator decomposition has no $\Delta S^z=0$ component creates a
state orthogonal to $\ket{\dd G}$ and cannot contribute to the matched
filter.

In the Cartesian Pauli basis, each $S^x$ or $S^y$ contributes one
transverse ladder operator:
\[
  S^x=\frac12(S^+ + S^-),
  \qquad
  S^y=\frac{1}{2i}(S^+ - S^-).
\]
A product with an odd number of transverse factors has only odd total
$\Delta S^z$ sectors and therefore no $\Delta S^z=0$ component.  Hence
we must also require
\[
  \#\{i:c_i\in\{1,2\}\}\equiv0\pmod 2.
\]

Therefore, for this real-state, $S^z_\pi$ problem, the code enumerates
only connected Pauli strings satisfying both
\[
  \#Y\equiv1\pmod2,
  \qquad
  \#(X,Y)\equiv0\pmod2.
\]
The first condition is the real-form/time-reversal filter; the second is
the $U(1)$ spin-charge filter.  Both are exact for this campaign.  As a
sanity check, on the dimer state the reduced basis gives
\[
  \eta_{\le2}=1,
  \qquad
  \eta_{\le4}=1,
\]
matching the full ED result.

\section{Operator Enumeration}

The local Pauli alphabet is encoded as
\[
  0=I,\qquad 1=S^x,\qquad 2=S^y,\qquad 3=S^z.
\]
For a window of length $w$, a code tuple
\[
  c=(c_1,\ldots,c_w)
\]
is included if
\begin{enumerate}
  \item it has exact support on that window:
  \[
    c_1\ne0,\qquad c_w\ne0,
  \]
  \item and it has odd $Y$ parity:
  \[
    \#\{i:c_i=2\}\equiv1\pmod 2.
  \]
  \item and it lies in the $\Delta S^z=0$ sector:
  \[
    \#\{i:c_i=1\ \text{or}\ c_i=2\}\equiv0\pmod 2.
  \]
\end{enumerate}
For a chain with open boundaries, windows are
\[
  W_{r,w}=\{r,r+1,\ldots,r+w-1\},
  \qquad
  r=1,\ldots,N-w+1.
\]
The cumulative basis for $\eta_{\le k}$ is
\[
  \mathcal B_{\le k}
  =
  \bigcup_{w=1}^{k}
  \bigcup_{r=1}^{N-w+1}
  \left\{
    O_{r,w,c}
  \right\}.
\]
Thus $\eta_{\le4}$ contains all $\eta_{\le2}$ operators plus the
additional three- and four-site connected-window strings.

\section{MPS-Level Algorithm}

For each saved checkpoint
\[
  \texttt{obs/T0\_path/<label>\_N<N>\_psi.h5},
\]
the code performs the following steps.

\subsection{Step 1: Load and Densify the MPS}

The production DMRG used $S^z$ quantum-number conservation.  Operators
containing $S^x$ and $S^y$ change quantum number sectors, so before
applying arbitrary Pauli strings we convert the block-sparse MPS to a
dense MPS:
\begin{lstlisting}
f = h5open(psi_path, "r")
psi = read(f, "psi", MPS)
close(f)

psi = dense(psi)
sites = siteinds(psi)
\end{lstlisting}

At the tensor level, an open-boundary MPS is
\begin{equation}
  \ket{\psi}
  =
  \sum_{\sigma_1,\ldots,\sigma_N}
  A^{[1]\sigma_1}
  A^{[2]\sigma_2}
  \cdots
  A^{[N]\sigma_N}
  \ket{\sigma_1\cdots\sigma_N}.
\end{equation}
Each local tensor has one physical index $\sigma_j$ and up to two bond
indices:
\begin{equation}
  A^{[j]}
  =
  A^{[j]\sigma_j}_{\alpha_{j-1},\alpha_j}.
\end{equation}
With quantum-number conservation turned on, the tensor is block sparse:
only entries obeying the $S^z$ flux rule are stored.  A local $S^x$ or
$S^y$ operator changes the $S^z$ quantum number on the physical index,
so applying it generally creates sectors that were absent in the
block-sparse representation.  The call
\[
  \texttt{dense(psi)}
\]
removes the QN block structure and stores ordinary dense tensors, so
non-$S^z$-conserving Pauli strings can be applied safely.

\subsection{Step 2: Build the Target Vector}

Construct the MPO
\[
  G=\sum_{j=1}^{N}(-1)^j S^z_j
\]
and apply it to the MPS:
\begin{lstlisting}
function staggered_Gpsi(psi::MPS, sites)
    N = length(psi)
    os = OpSum()
    for j in 1:N
        os += (-1)^j, "Sz", j
    end
    G = MPO(os, sites)
    out = apply(G, psi; cutoff=1e-14)
    noprime!(out)
    return out
end
\end{lstlisting}
Then center the vector:
\begin{lstlisting}
function centered(phi::MPS, psi::MPS)
    ev = inner(psi, phi)
    return phi - ev * psi
end

target = centered(staggered_Gpsi(psi, sites), psi)
\end{lstlisting}
Mathematically,
\[
  \ket{\mathrm{target}}
  =
  \ket{\dd G}
  =
  \left(G-\langle G\rangle\right)\ket{\psi}.
\]

\subsection{Step 3: Apply Each Pauli String}

For each connected-window Pauli string $O_a$, apply the local operators
directly to a copy of the MPS:
\begin{lstlisting}
const PAULI_OPS = Dict(1 => "Sx", 2 => "Sy", 3 => "Sz")

function apply_pauli_string(psi::MPS, sites, start::Int, codes::Vector{Int})
    phi = copy(psi)
    for (off, code) in enumerate(codes)
        code == 0 && continue
        j = start + off - 1
        Oj = dense(op(PAULI_OPS[code], sites[j]))
        phi[j] = Oj * phi[j]
        noprime!(phi[j])
    end
    return phi
end
\end{lstlisting}
Then center it:
\begin{lstlisting}
phi_a = centered(apply_pauli_string(psi, sites, start, codes), psi)
\end{lstlisting}
This gives
\[
  \ket{\phi_a}
  =
  \left(O_a-\langle O_a\rangle\right)\ket{\psi}.
\]

\subsection{What local operator application means}

Suppose a one-site operator $O_j$ has matrix elements
\[
  (O_j)_{\tau_j\sigma_j}
\]
on the physical spin-$1/2$ basis.  Applying it to the MPS changes only
the local tensor at site $j$:
\begin{equation}
  \widetilde A^{[j]\tau_j}_{\alpha_{j-1},\alpha_j}
  =
  \sum_{\sigma_j}
  (O_j)_{\tau_j\sigma_j}
  A^{[j]\sigma_j}_{\alpha_{j-1},\alpha_j}.
  \label{eq:local_tensor_application}
\end{equation}
All other tensors are unchanged.  A Pauli string on a connected window
\[
  O_a
  =
  P_{r}^{(c_1)}
  P_{r+1}^{(c_2)}
  \cdots
  P_{r+w-1}^{(c_w)}
\]
is applied by repeating \eqref{eq:local_tensor_application} on each site
in the window.  This is precisely what the code
\begin{lstlisting}
phi[j] = Oj * phi[j]
noprime!(phi[j])
\end{lstlisting}
does.  In ITensor notation, \texttt{Oj} carries a primed output site
index and an unprimed input site index.  Multiplying it into
\texttt{phi[j]} contracts the input physical leg with the tensor's site
index and leaves a primed output site index.  The call
\texttt{noprime!} renames that output leg back to the standard site
index so the resulting object is again a valid MPS with the same site
index set as $\ket{\psi}$.

No global Hilbert-space vector is formed.  The cost of applying a
length-$w$ Pauli string is just $w$ local tensor contractions, followed
by MPS inner products when that vector is used in the Gram matrix.

\subsection{Step 4: Build the Gram Matrix and Target Vector}

The code never forms a $2^N$ Hilbert-space vector.  Instead it evaluates
all needed quantities with MPS inner products:
\begin{lstlisting}
K[a,c] = real(inner(phis[a], phis[c]))
b[a]   = imag(inner(phis[a], target))
\end{lstlisting}
The real symmetric matrix
\[
  K_{ab}
  =
  \mathrm{Re}\braket{\phi_a}{\phi_b}
\]
is the covariance/Gram matrix of the real operator-action subspace.
The vector
\[
  b_a
  =
  \mathrm{Im}\braket{\phi_a}{\dd G}
\]
is the real-form target overlap.  Up to the common normalization
conventions, this is the matched-filter signal $s_a$.

\subsection{How MPS inner products are evaluated}

For two MPSs
\[
  \ket{\phi}
  =
  \sum_{\bm\sigma}
  B^{[1]\sigma_1}\cdots B^{[N]\sigma_N}
  \ket{\bm\sigma},
  \qquad
  \ket{\chi}
  =
  \sum_{\bm\sigma}
  C^{[1]\sigma_1}\cdots C^{[N]\sigma_N}
  \ket{\bm\sigma},
\]
the inner product is the tensor-network contraction
\begin{equation}
  \braket{\phi}{\chi}
  =
  \sum_{\bm\sigma}
  \overline{
    B^{[1]\sigma_1}
    \cdots
    B^{[N]\sigma_N}
  }
  C^{[1]\sigma_1}
  \cdots
  C^{[N]\sigma_N}.
\end{equation}
ITensor evaluates this from left to right by building an environment:
\begin{equation}
  L_j
  =
  \sum_{\sigma_1,\ldots,\sigma_j}
  \overline{
    B^{[1]\sigma_1}\cdots B^{[j]\sigma_j}
  }
  C^{[1]\sigma_1}\cdots C^{[j]\sigma_j}.
\end{equation}
This avoids ever expanding the $2^N$ coefficient vector.  The final
scalar is what \texttt{inner(phi, chi)} returns.

Thus the Gram element
\[
  K_{ab}=\mathrm{Re}\braket{\phi_a}{\phi_b}
\]
is computed by contracting two MPSs produced by local Pauli-string
actions.  Similarly,
\[
  b_a=\mathrm{Im}\braket{\phi_a}{\dd G}
\]
is computed by contracting the Pauli-string-action MPS with the MPO
application MPS $\ket{\dd G}$.

\subsection{Step 5: Solve the Projection / Wiener Problem}

Diagonalize
\[
  K=V\Lambda V^T.
\]
Discard eigenvalues below a numerical tolerance
\[
  \lambda_\alpha
  \le
  \epsilon_{\rm rel}\max(1,\lambda_{\rm max}).
\]
Then compute
\begin{equation}
  \mathrm{num}
  =
  \sum_{\alpha:\lambda_\alpha>0}
  \frac{
    \left[(V^Tb)_\alpha\right]^2
  }{
    \lambda_\alpha
  }.
\end{equation}
Finally
\begin{equation}
  \boxed{
  \eta_{\le k}
  =
  \frac{\mathrm{num}}{\braket{\dd G}{\dd G}}.
  }
\end{equation}
The corresponding code is:
\begin{lstlisting}
function eta_from_basis(phis::Vector{MPS}, target::MPS; eig_rtol::Float64=1e-10)
    M = length(phis)
    K = zeros(Float64, M, M)
    b = zeros(Float64, M)

    for a in 1:M
        b[a] = imag(inner(phis[a], target))
        for c in 1:a
            val = real(inner(phis[a], phis[c]))
            K[a, c] = val
            K[c, a] = val
        end
    end

    K = 0.5 .* (K .+ K')
    evals, evecs = eigen(Symmetric(K))

    maxeval = maximum(abs.(evals))
    tol = max(1.0, maxeval) * eig_rtol
    keep = findall(>(tol), evals)

    coeff = evecs[:, keep]' * b
    num = sum(abs2.(coeff ./ sqrt.(evals[keep])))
    norm_target = real(inner(target, target))

    return num / norm_target, length(keep)
end
\end{lstlisting}

\section{Cumulative Computation of \texorpdfstring{$\eta_{\le2}$}{eta<=2}
and \texorpdfstring{$\eta_{\le4}$}{eta<=4}}

The implementation accumulates the basis by increasing window size.  For
each $k$, all strings of length $k$ are appended to the existing basis,
and the projection is recomputed:
\begin{lstlisting}
phis = MPS[]

for k in 1:kmax
    for start in 1:(length(psi) - k + 1)
        for codes in odd_y_codes(k)
            phi = centered(apply_pauli_string(psi, sites, start, codes), psi)
            push!(phis, phi)
        end
    end

    eta, rank = eta_from_basis(phis, target; eig_rtol=eig_rtol)
    eta_by_k[string(k)] = eta
    rank_by_k[string(k)] = rank
    nops_by_k[string(k)] = length(phis)
end
\end{lstlisting}
Thus the reported $k=4$ value is truly
\[
  \eta_{\le4},
\]
not the contribution of exactly four-site windows alone.

\section{Submitted Production Job}

The current production postprocessing is submitted as
\[
  \texttt{qfi\_eta\_path\_T0.slurm}.
\]
It runs one path checkpoint per array task:
\begin{lstlisting}[language=bash]
julia --project="$QFI_JULIA_PROJECT" compute_eta_path_T0.jl \
    --input-file "$IN" \
    --out "$OUT" \
    --kmax 4 \
    --eig-rtol 1e-10
\end{lstlisting}
The output for each state is
\[
  \texttt{obs/T0\_path\_eta/<label>\_N<N>\_eta.json}.
\]
Each JSON contains
\begin{itemize}
  \item the path label and couplings $(J_1,J_2,J_3)$,
  \item $N$, $E_0$, $\chi_{\rm final}$, and $\FQ$,
  \item $\eta_{\le1},\eta_{\le2},\eta_{\le3},\eta_{\le4}$,
  \item the numerical rank retained in each Gram solve,
  \item the number of operator-action vectors included.
\end{itemize}

\section{Sanity Check}

Before launching the full path job, the code was tested on the corrected
dimer state $D=(1,0,0)$ at $N=16$.  Since this is a product of isolated
singlets, the matched filter must saturate already at $k=2$:
\[
  \eta_{\le2}(D)=1.
\]
The MPS postprocessor returned
\[
  \eta_{\le2}=1.0000000000000016,
  \qquad
  \eta_{\le4}=1.0000000000000007.
\]
This confirms the operator basis, real-form projection, and MPS inner
product implementation are consistent with the ED matched-filter
definition.

\section{Final Figure to Produce}

After all 39 postprocessing tasks complete, the final important figure is
\[
  \eta_{\le2},\ \eta_{\le4}
  \quad\text{along}\quad
  D\to C\to U\to D
\]
for
\[
  N=16,32,64.
\]
Together with the already-generated $F_Q$ path figure, this gives the two
central numerical diagnostics for the chosen generator and chosen
observable algebra:
\[
  \boxed{\FQ\ \text{and}\ \eta_{\le k}.}
\]

The interpretation is:
\begin{itemize}
  \item $\FQ$ gives the total staggered response strength.
  \item $\eta_{\le k}$ gives the fraction of that response visible to
  connected local windows of size at most $k$.
\end{itemize}
Therefore the pair
\[
  \left(\FQ,\eta_{\le k}\right)
\]
separates ``how much response exists'' from ``how that response can be
read out.''  Along the corrected
$D\to C\to U\to D$ path, the expectation is:
\begin{itemize}
  \item at the dimer point $D$, the response is visible already through
  two-site bond operators;
  \item at the four-site cluster point $C$, the response should be
  visible through four-site windows;
  \item at the uniform Heisenberg point $U$, fixed finite windows should
  capture a decreasing fraction as $N$ grows, because the response is
  scale-delocalized.
\end{itemize}
This statement is intentionally about the \emph{operator channel} of the
response.  In another phase, such as a Haldane SPT phase, the correct
operator channel may be string-like rather than local-window-like.  In
that case one should compare
\[
  \eta_{\rm local}(k)
  \qquad\text{against}\qquad
  \eta_{\rm string}(k,\ell),
\]
rather than interpreting a small local-window value as long-range
entanglement.

\end{document}
